\documentclass[../../../include/open-logic-section]{subfiles}

\begin{document}

\olfileid{sth}{z}{story}
\olsection{د کيسې نور تفصيل}
	
په \olref[story][approach]{sec} کښې مو د سټونو د جوړېدو د بهير په اړه د
شوينفيلډ بيان رانقل کړے و. اوس غواړو يو څو نور اصول وليکو، څو دا
کيسه لږه نوره دقيقه شي. اصول دا دي:
\begin{enumerate}
\item[] \stageshier. هر سټ په يوه پړاو کښې جوړېږي.
\item[] \stagesord. پړاوونه ترتيب لري: ځينې د نورو \emph{نه مخکې}
راځي.\footnote{په حقيقت کښې به په پټه دا فرض کوو چې پړاوونه
\emph{ښه ترتيب شوي} دي. د دې مانا په \olref[ordinals][]{chap}
کښې تشريح کېږي. دا يو مهم فرض دے. په حقيقت کښې، د \citet{Scott1974}
د يوې ډېرې ځيرکې طريقې په کارولو سره، له دې فرضه \emph{ډډه}
کېدلے شي او بيا ترې \emph{نتيجه} راايستل کېدلے شي. (دا به دا هم
تشريح کړي چې ولې بايد يو لومړنے پړاو ومنو.) دلته دې بحث ته نۀ شو
ننوتلے؛ د نورو معلوماتو دپاره \citet{ButtonLT1} وګورئ.}
\item[] \stagesacc. د هر پړاو $S$ او د هغو سټونو د هرې ټولګې
دپاره چې له پړاو~$S$ \emph{نه مخکې} جوړ شوي وي: په پړاو~$S$
کښې يو سټ جوړېږي چې غړي ئې په کره توګه هماغه سټونه وي.
په پړاو~$S$ کښې بل هېڅ شے نۀ جوړېږي.
\end{enumerate}
دا غير رسمي اصول دي، خو د زرمېلو د سټ تيورۍ د څو بديهي اصلونو
د توجيه دپاره به ئې وکارولے شو.

(دلته يو احتياط بيانول پکار دي. که څه هم هغه بديهي اصلونه چې
وړاندې کوو ئې، او د هغو نومونه، په بشپړه توګه معياري دي، هغه په
مايل ليک ليکل شوي اصول چې اوس مو وړاندې کړل په اړوندو ليکنو کښې
ځانګړي نومونه نۀ لري. موږ ورته يوازې داسې نومونه غوره کړي دي
چې هيله لرو ګټور به وي.)

\end{document}
